\chapter{Mise en œuvre de l'ETL avec Talend}
\label{ch:etl}

\section{Introduction}
Ce chapitre détaille la mise en œuvre du processus ETL (Extract-Transform-Load) permettant d'alimenter l'entrepôt de données à partir des sources brutes (fichiers Excel/CSV issus de l'EMD et des comptages de trafic, préalablement nettoyés sous Python/Colab). L'outil Talend Open Studio a été utilisé pour concevoir cette chaîne sous forme de \emph{jobs} graphiques.

\section{Architecture générale de l'ETL}
La chaîne ETL est organisée en deux niveaux successifs~:
\begin{enumerate}
\item une \textbf{zone de préparation (Staging Area, SA)}~: jobs chargeant les données brutes nettoyées dans des tables intermédiaires, sans transformation lourde~;
\item une \textbf{zone de chargement final}~: jobs lisant les tables de la Staging Area, appliquant les transformations, jointures et calculs nécessaires (via le composant \texttt{tMap}), puis chargeant les dimensions et les faits dans l'entrepôt de données (DW).
\end{enumerate}

\section{Zone de préparation (Staging Area)}
Trois jobs principaux assurent l'alimentation de la zone de préparation~: \texttt{SA\_individu}, \texttt{SA\_menage} et \texttt{SA\_deplacement}, correspondant aux trois grandes entités sources de l'EMD (individus, ménages, déplacements). Chaque job suit un schéma simple~: lecture du fichier source, filtrage des colonnes utiles via \texttt{tFilterColumns}, puis chargement dans la base via \texttt{tDBOutput}.

\begin{figure}[H]
\centering
\includegraphics[width=0.85\textwidth]{talend/SA_deplacement.png}
\caption{Job Talend de chargement de la Staging Area -- déplacements}
\label{fig:sa_deplacement}
\end{figure}

\section{Configuration des dimensions}
La configuration des dimensions géographiques et de site illustre la logique de mapping utilisée pour la construction des tables de référence de l'entrepôt.

\begin{figure}[H]
\centering
\begin{subfigure}[b]{0.48\textwidth}
\includegraphics[width=\textwidth]{talend/config_dim_site.png}
\caption{Configuration de \texttt{Dim\_Site}}
\end{subfigure}
\hfill
\begin{subfigure}[b]{0.48\textwidth}
\includegraphics[width=\textwidth]{talend/dim_geo.png}
\caption{Chargement de \texttt{Dim\_Geographie}}
\end{subfigure}
\caption{Jobs de construction des dimensions de référence}
\end{figure}

De même, les dimensions \texttt{Dim\_Individu} et \texttt{Dim\_Menage} sont alimentées par des jobs dédiés, reprenant les données nettoyées issues de la Staging Area correspondante.

\section{Construction des tables de faits}

\subsection{Fait\_Deplacement}
Le job \texttt{Job\_deplacement} illustre la logique de transformation appliquée pour construire la table de faits \texttt{Fact\_deplacement}. Le composant \texttt{tMap} central réalise les opérations suivantes (figure~\ref{fig:tmap_depl})~:
\begin{itemize}
\item récupération du flux principal (\texttt{row1}) contenant les attributs du déplacement (numéro de ménage, numéro d'individu, lieux et heures de départ/arrivée, durée, motif, fréquence de transport en commun, mode, coût total)~;
\item trois jointures (\emph{lookup}) sur les dimensions \texttt{Dim\_Motif} (par \texttt{motif}), \texttt{Dim\_Individu} (par clé composite \texttt{num\_individu + num\_menage}) et \texttt{Dim\_Mode} (par \texttt{mode}) afin de résoudre les clés étrangères~;
\item construction du flux de sortie \texttt{fact\_deplacement} combinant les clés étrangères résolues (\texttt{fk\_id\_motif}, \texttt{fk\_id\_individu}, \texttt{fk\_id\_transport}) et les mesures (durée du trajet, lieux, coût total, heures de départ/arrivée, fréquence de transport en commun).
\end{itemize}

\begin{figure}[H]
\centering
\includegraphics[width=\textwidth]{talend/tmap_depl.png}
\caption{Composant tMap du job de construction de \texttt{Fact\_deplacement}}
\label{fig:tmap_depl}
\end{figure}

\subsection{Fait\_Comptage}
Le job de construction de \texttt{Fact\_Comptage} suit une logique similaire~: une entrée principale (\texttt{tDBInput\_1}) portant les mesures de comptage est enrichie par deux jointures de référence (\texttt{tDBInput\_2}, \texttt{tDBInput\_3}) via un composant \texttt{tMap\_1}, avant chargement final dans la table \texttt{fact\_comptage} de l'entrepôt.

\begin{figure}[H]
\centering
\includegraphics[width=0.85\textwidth]{talend/fact_compt.png}
\caption{Job Talend de construction de \texttt{Fact\_Comptage}}
\label{fig:fact_compt}
\end{figure}

\subsection{Fait\_Accessibilité et Fait\_IndividuMenage}
Les jobs \texttt{fact\_accessibilite} et \texttt{fact\_indivmen} suivent le même schéma général, en exploitant respectivement les variables d'accessibilité (durée de marche, modes de transport disponibles, exposition aux inondations) et les variables croisées individu/ménage (taille du ménage, équipement, motorisation).

\section{Composants Talend mobilisés}
Le tableau~\ref{tab:talend_comp} récapitule les principaux composants Talend utilisés dans les différents jobs.

\begin{table}[H]
\centering
\caption{Composants Talend utilisés}
\label{tab:talend_comp}
\begin{tabular}{ll}
\toprule
\textbf{Composant} & \textbf{Rôle} \\
\midrule
tFileInputExcel / tFileInputDelimited & Lecture des fichiers sources (EMD, comptages) \\
tFilterColumns & Sélection des colonnes utiles (priorisation MVP) \\
tMap & Jointures, transformations et calculs de mesures \\
tDBInput & Lecture de tables intermédiaires / dimensions \\
tDBOutput & Chargement dans la Staging Area ou l'entrepôt \\
\bottomrule
\end{tabular}
\end{table}

\section{Conclusion}
Le processus ETL conçu avec Talend assure une alimentation fiable et traçable de l'entrepôt de données, en respectant la priorisation des données définie au chapitre~\ref{ch:analyse} et le modèle décisionnel conçu au chapitre~\ref{ch:modele}. Le chapitre suivant présente l'implémentation physique de cet entrepôt sous PostgreSQL.

\clearpage
